\begin{onehalfspace}

\section{Discussione complessiva dei risultati}
\label{sec:discussione_complessiva}

I risultati riguardano tre modi di intervenire: cercare candidati nell'istogramma finale, proteggere lo stato con la QEC e ridurre le porte del circuito. Ogni intervento ha una metrica diversa. Nella campagna NISQ si misurano le iterazioni; nei test QEC, i fallimenti logici; nell'analisi di Shor, la sensibilità a un errore Pauli per gate. Queste quantità non si convertono direttamente fra loro. Permettono però di capire quale problema risolva ciascun metodo.

Nel regime NISQ il risultato più netto non appartiene al classificatore, ma alla regola di
ricerca. Nei due scenari appaiati, TOP-4 porta il numero medio di iterazioni da $1.37$ e
$1.50$ a $1.00$, con successo in tutte le repliche, mentre M2 non migliora significativamente
la baseline e risulta peggiore della propria ablazione priva di filtro
(Tabella~\ref{tab:riepilogo_rho}). Anche la ZNE digitale non riduce le iterazioni e porta gli
shot medi totali a $3\,072$ e $4\,813$, contro i $1\,024$ di TOP-4
(Tabella~\ref{tab:confronto_zne}). In questo perimetro, dunque, conviene prima sfruttare in
modo esplicito i pochi picchi ancora informativi; un modello appreso a valle non ricostruisce
informazione che il rumore ha già cancellato e deve essere valutato contro una regola semplice
disponibile sul medesimo istogramma.

I laboratori QEC mostrano invece che intervenire durante la conservazione dello stato può
cambiare l'ordine dell'errore, purché si operi nel regime appropriato. Il codice a ripetizione
segue $p_L=3p^2-2p^3$ entro l'incertezza Monte Carlo
(Figura~\ref{fig:qec_repetition}); lo Steane $[[7,1,3]]$ estende la protezione a errori
arbitrari, con pendenza log--log $1.94$ e pseudo-soglia attorno a $p\simeq0.08$
(Figura~\ref{fig:qec_steane}). Il surface code rende scalabile la stessa idea: i fit forniscono
$p_{th}=0.86\%$ in base Z e $0.81\%$ in base X e, a $p=2\times10^{-3}$, l'errore logico
scende da $1.94\times10^{-3}$ a distanza 3 a $2.01\times10^{-5}$ a distanza 9
(Figura~\ref{fig:qec_surface}). Questi valori costituiscono una baseline entro il modello
Pauli circuit-level; non sono ancora una previsione calibrata di un'esecuzione
fault-tolerant di Shor.

La campagna sulla decodifica chiarisce quando l'apprendimento aiuta. Sostituire MWPM con una rete non dà vittorie nelle quattordici configurazioni provate. Per pareggiare a distanza 3 servono $10^5$--$10^6$ campioni. Quando il rumore contiene correlazioni che il grafo del matching non rappresenta, l'ibrido riduce invece l'errore fino al $21\%$. Il vantaggio diminuisce a distanza 5 e si annulla a distanza 7 (Tabella~\ref{tab:sintesi_decodifica}). A distanza 5, BP+OSD recupera già il $17$--$27\%$, contro l'$1$--$3\%$ dell'ibrido. Il modello appreso aiuta quindi a riconoscere correlazioni assenti dal metodo di riferimento. Una calibrazione imprecisa, da sola, non spiega il beneficio.

L'analisi di sensibilità di Shor collega questi risultati senza forzare una conversione tra
metriche. Nel proxy uniforme per gate, il successo per singola misura passa da $0.7489$ in
assenza di rumore a $0.4518$ per $p_L=0.02$, quindi raggiunge un plateau empirico vicino a
$0.245$ oltre $p_L=0.1$ (Figura~\ref{fig:qec_shor_logico}). TOP-4 resta pari a $1.000$ fino
a $p_L=0.05$, ma scende a $0.670$ a $p_L=0.2$: il post-processing conserva quindi efficacia
anche quando la massa utile è degradata, senza rendere il circuito immune al rumore. Come
formalizzato nella Tabella~\ref{tab:decoder_su_shor}, trasformare le curve QEC in probabilità
di fattorizzazione richiederebbe ancora un gate set logico, cicli di sindrome, operazioni
non-Clifford e un modello delle correlazioni residue.

Infine, la QFT approssimata è utile quando il rumore evitato compensa la precisione persa. Su Shor per $N=21$, le porte a due qubit scendono da $49\,661$ a $23\,178$, ma il successo ideale cala da $0.4667$ a $0.4013$. I circa tre punti percentuali di possibile vantaggio sono una stima del modello ausiliario, che richiede errori a due qubit circa duecento volte inferiori a quelli della calibrazione considerata. Nei test rumorosi questo beneficio non è osservato. Sulla sola stima di fase, la variante troncata ha invece un risultato migliore in tutti e nove i confronti e raggiunge $+0.126$ su dieci qubit. Il grado migliore resta due o tre alle taglie provate (Sezione~\ref{sec:qft_topologia}). Il costo del circuito determina quindi se la stessa approssimazione sia conveniente.

Nel complesso emerge una regola unitaria: la complessità è giustificata solo quando agisce sul
collo di bottiglia osservato. TOP-K sfrutta informazione residua; la QEC sopprime l'errore sotto
soglia; l'apprendimento recupera struttura assente dal decoder analitico; il troncamento della
QFT conviene se il risparmio circuitale supera la perdita algoritmica. Restano tre limiti:
l'istanza rumorosa principale è $N=15$, le soglie dipendono dal modello e manca una simulazione
end-to-end fault-tolerant. Sono quindi criteri sperimentali, non previsioni per QPU reali.

\end{onehalfspace}
